\subsection{Циклы}

Циклы позволяют выполнять один и тот же блок кода многократно.

\subsubsection*{Цикл for}
Используется, когда известно количество повторений. Функция \texttt{range(N)} создает последовательность от 0 до N-1.

\subsubsection*{Цикл while}
Выполняется, пока условие истинно. Будьте осторожны: если условие никогда не станет ложным, получится <<бесконечный цикл>>.

\begin{codefile}[python]{python/01_python_basics/06_loops.py}
\end{codefile}

\begin{itemize}
    \item \texttt{import time} подключает стандартный модуль \texttt{time}; \texttt{time.sleep(0.5)} --- пауза на полсекунды. В отличие от Lua на борту, Python-программа работает на компьютере (в симуляторе --- в браузере) и автопилот не блокирует, поэтому \texttt{sleep} здесь допустим (глава~\ref{sec:python-sdk}).
    \item \texttt{f"Мигание №\{i+1\}"} --- f-строка: выражение в фигурных скобках вычисляется и подставляется в текст. \texttt{battery -= 2} --- то же, что \texttt{battery = battery - 2}.
    \item \texttt{break} — немедленно прерывает цикл.
    \item \texttt{continue} — переходит к следующей итерации.
\end{itemize}

\subsubsection*{Список и цикл по списку}
Несколько значений в определённом порядке хранят в \textbf{списке} --- значения через запятую в квадратных скобках. Цикл \texttt{for} перебирает элементы списка по очереди, поэтому одно и то же действие применяется ко всем:

\begin{codebox}[Список сигналов]
\begin{verbatim}
signals = ["ожидание", "готовность", "завершение"]
for signal in signals:
    print(signal)
print("Всего:", len(signals), "первый:", signals[0])
\end{verbatim}
\end{codebox}

Нумерация элементов в Python начинается с нуля: \texttt{signals[0]} --- первый элемент, \texttt{signals[2]} --- третий. Функция \texttt{len()} возвращает длину списка. В главе~\ref{sec:python-sdk} списком будут вершины маршрута.

\subsubsection*{Задачи для самостоятельного решения}

\begin{enumerate}
    \item Выведите числа от 0 до 4 используя \texttt{range(5)}.
    \item Выведите слово "Привет" 3 раза.
    \item Выведите числа от 1 до 5 (используйте \texttt{range(1, 6)}).
    \item Выведите четные числа от 0 до 8 (\texttt{range(0, 9, 2)}).
    \item Найдите сумму чисел от 1 до 100.
    \item Выведите числа от 5 до 1 (\texttt{range(5, 0, -1)}).
    \item Выведите таблицу умножения на 5 (от 1*5 до 10*5).
    \item Дана строка "Дрон". Выведите каждую букву с новой строки.
    \item Используя \texttt{while}, выводите \texttt{x}, пока \texttt{x < 3}. Начальное \texttt{x = 0}.
    \item Ищите число 7 в диапазоне от 0 до 10. Если нашли — используйте \texttt{break}.
    \item Выведите числа от 0 до 4, пропустив 2 (используйте \texttt{continue}).
    \item Выведите квадраты чисел от 1 до 5.
    \item Посчитайте длину строки "Hello" с помощью цикла (без \texttt{len()}).
    \item Выведите 5 строк, в каждой по одной звездочке \texttt{*}.
    \item Выведите лесенку из звезд (1, 2, 3...).
\end{enumerate}

\begin{taskbox}[Обратный отсчёт]
С помощью \texttt{range} с отрицательным шагом выведите числа от 10 до 1, каждое с новой строки, а затем --- \texttt{Поехали!}. Цикл должен выполниться ровно 10 раз.
\expected{11 строк: \texttt{10}, \texttt{9}, \ldots, \texttt{1}, \texttt{Поехали!}.}
\mission{Python}{24}
\end{taskbox}

\begin{taskbox}[Список сигналов]
Создайте список из трёх названий сигналов: «ожидание», «готовность», «завершение». Выведите каждый элемент циклом \texttt{for}, а затем количество элементов через \texttt{len()}. Добавьте в список четвёртый сигнал и убедитесь, что цикл менять не пришлось.
\expected{три (после правки --- четыре) строки с названиями и строка \texttt{Всего сигналов: 3} (затем \texttt{4}). Решения обоих заданий --- в конце файла \texttt{examples/python-answers/06\_loops\_answers.py}.}
\mission{Python}{27}
\end{taskbox}
